\documentclass[aspectratio=169,10pt,english]{beamer}

\input{../../../_preambulo_beamer.tex}
\input{../../../_comandos_beamer.tex}

\selectlanguage{english}

\title{MA1001B - Statistical Modeling for Decision Making}
\subtitle{Lesson 09: Conditional Probability and Independence}
\author{}
\institute{Tecnol\'ogico de Monterrey}
\date{}

\begin{document}

\begin{frame}[plain]
  \vspace{1.2cm}
  {\Large\bfseries MA1001B - Statistical Modeling for Decision Making\par}
  \vspace{0.7cm}
  {\large Lesson 09: Conditional Probability and Independence\par}
  \vspace{0.7cm}
  {\color{TecRojo}\rule{\textwidth}{1.2pt}}
  \vspace{0.8cm}

  MA1001B Faculty\\[0.3cm]
  Term\\[0.3cm]
  Tecnol\'ogico de Monterrey
\end{frame}

\begin{frame}{The Conditional Probability Question}
  \begin{alertblock}{Guiding question}
    How does knowing one event occurred change the relevant sample space, and
    when does that knowledge leave another probability unchanged?
  \end{alertblock}

  \vspace{0.4em}
  By the end of this lesson, you can:
  \begin{itemize}
    \item calculate and interpret conditional probabilities;
    \item explain why reversing a condition changes the question;
    \item test independence with equivalent probability equalities;
    \item update a sequence probability when sampling without replacement.
  \end{itemize}

  \vfill
  \scriptsize All order records and operating processes are synthetic.
\end{frame}

\begin{frame}{Conditioning Reduces the Sample Space}
  \small
  Let $M$ be manual review and $L$ be late delivery.

  \begin{center}
    \begin{tabular}{lrrr}
      \toprule
      \textbf{Routing process} & \textbf{Late} & \textbf{On time} & \textbf{Total} \\
      \midrule
      Manual review & 48 & 72 & 120 \\
      Automated & 32 & 248 & 280 \\
      \midrule
      Total & 80 & 320 & 400 \\
      \bottomrule
    \end{tabular}
  \end{center}

  \[
    P(L\mid M)=\frac{P(L\cap M)}{P(M)}=\frac{48}{120}=0.40.
  \]

  The condition $M$ makes the 120 manually reviewed orders the new reference
  group, rather than all 400 orders.
\end{frame}

\begin{frame}{Step 1: Direction Changes the Denominator}
  \begin{columns}[T]
    \begin{column}{0.36\textwidth}
      \small
      \begin{block}{Verified output}
        $P(L\mid M)=0.400000$\\
        $P(L\mid A)=0.114286$\\
        $P(M\mid L)=0.600000$
      \end{block}

      $P(L\mid M)$ asks about late orders within manual review. $P(M\mid L)$
      asks about manual review within late orders. They need not match.
    \end{column}
    \begin{column}{0.64\textwidth}
      \centering
      \includegraphics[width=0.93\textwidth]{../figures/conditional_01_rates.png}
    \end{column}
  \end{columns}
\end{frame}

\begin{frame}{Independence Means No Probability Update}
  \small
  Events $M$ and $L$ are independent if any equivalent test holds:
  \[
    P(L\mid M)=P(L),
    \qquad
    P(M\cap L)=P(M)P(L).
  \]

  For the synthetic orders:
  \[
    0.40\ne0.20,
    \qquad
    0.12\ne(0.30)(0.20)=0.06.
  \]

  \begin{alertblock}{Interpretation limit}
    The table shows an association between routing and delivery status. It does
    not show that manual review causes late delivery.
  \end{alertblock}
\end{frame}

\begin{frame}{Step 2: Two Equivalent Tests Reach One Decision}
  \begin{columns}[T]
    \begin{column}{0.34\textwidth}
      \small
      \begin{block}{Verified comparison}
        Conditional: $0.40\ne0.20$\\
        Product: $0.12\ne0.06$\\[0.2em]
        Independent: \textbf{False}
      \end{block}

      Only one failed equality is needed to reject independence. The second is
      a consistency check.
    \end{column}
    \begin{column}{0.66\textwidth}
      \centering
      \includegraphics[width=0.98\textwidth]{../figures/conditional_02_independence.png}
    \end{column}
  \end{columns}
\end{frame}

\begin{frame}{Step 3: Without Replacement Creates Dependence}
  \begin{columns}[T]
    \begin{column}{0.39\textwidth}
      \small
      In a 10-order batch, 4 are flagged. For two selections without
      replacement,
      \[
        P(F_1\cap F_2)=\frac{4}{10}\frac{3}{9}=0.133333.
      \]

      Naive independence gives $0.4^2=0.160000$. A seed-42 simulation of
      500,000 selections gives \textbf{0.132938}.

      \textbf{Limit:} longer sequences require explicit branches.
    \end{column}
    \begin{column}{0.61\textwidth}
      \centering
      \includegraphics[width=\textwidth]{../figures/conditional_03_sequence.png}
    \end{column}
  \end{columns}
\end{frame}

\begin{frame}{Concept Check}
  Use the synthetic order table and inspection batch.
  \begin{enumerate}
    \item What denominator belongs in $P(L\mid M)$?
    \item Why is $P(L\mid M)=0.40$ different from $P(M\mid L)=0.60$?
    \item Are $M$ and $L$ independent? Support the decision with one
      equality test.
    \item After one flagged order is selected without replacement, why does the
      next flagged probability become $3/9$?
  \end{enumerate}

  \begin{block}{Connection to Lesson 10}
    Tree diagrams make changing conditional probabilities and multi-stage
    sample spaces visible, branch by branch.
  \end{block}

  \vfill
  \scriptsize Source: OpenStax, \textit{Introductory Business Statistics 2e},
  Chapter 3, Sections 3.1--3.3. CC BY-NC-SA 4.0.
\end{frame}

\appendix



\end{document}
